% Standalone insertion for the Two-Law technical synthesis.
% Assumes the macros and theorem environments of the main manuscript.

\section{Critical Boundary Actualization and Numerical Closure}
\label{sec:critical-boundary-actualization}

Primitive-Parameter Provenance says that a publicly inequivalent continuous coupling
cannot remain a bare dial.  Reversible Determination and Isometric Ledger Invariance
have now fixed the finite ordinary state and process arena, but they do not fix the
numerical generator of motion.  The remaining freedom sits chiefly in the effective
action and its vacuum: gauge kinetic coefficients, Yukawa eigenvalues, scalar
self-couplings, and the scale at which their boundary conditions are imposed.

The wider philosophy points to a specific next clause.  The Selective Law refuses an
unrecorded preference between complete payments; the Constitutive Law refuses an
unwitnessed distinction; and Public Confluence requires the stable public account to
be compatible across contexts.  Applied to phases rather than presentations, the
common content is:
\begin{center}
\emph{When complete phases satisfy the same inherited obligations and no prior
public record distinguishes them, the fundamental law may not hide a phase bias.
Their distinction must be actualized at a common stationary boundary.}
\end{center}

\begin{definition}[Unbiased public phases]
Let $\Gamma_{\mathrm{eff}}[\phi;\mathbf c]$ be the complete renormalized public
effective action of a selected physical API, with dimensionless couplings
$\mathbf c$.  An \emph{unbiased public phase family} is a finite family of
homogeneous configurations $\phi_i$ such that:
\begin{enumerate}[label=\textup{B\arabic*},leftmargin=2.4em]
\item each $\phi_i$ is a constitutively admissible stable or marginally stable
      phase of the same API;
\item every member discharges the same frozen structural and public-record
      obligations;
\item no inherited charge, boundary condition, topological sector, source,
      environment label, or locked historical record privileges one member; and
\item the comparison of their vacuum functionals is itself a lawful public record.
\end{enumerate}
The clauses are fixed before the numerical values of $\mathbf c$ are inspected.
A phase is not allowed to issue the record that retrospectively licenses its own
preference.
\end{definition}

\begin{axiom}[Critical Boundary Actualization]
\label{ax:critical-boundary-actualization}
For every unbiased public phase family in a completed fundamental realization,
there is an endogenously certified matching scale $\mu_*$ at which the phases are
stationary and have one common public vacuum functional:
\[
 \left.\frac{\delta\Gamma_{\mathrm{eff}}}{\delta\phi}\right|_{\phi_i,\mu_*}=0,
 \qquad
 \mathcal E_{\mathrm{vac}}(\phi_i;\mu_*)
 =\mathcal E_{\mathrm{vac}}(\phi_j;\mu_*)
 \quad\text{for all }i,j.
\]
The scale $\mu_*$ is not a new dial: it must be inherited from a locked record,
fixed by a spectral or gravitational interface, or generated as an RG-invariant
transmutation scale.  If no such placement certificate exists, the criticality
condition is a relation without a completed physical location.
\end{axiom}

\begin{remark}[No optimizer and no universal claim of metastability]
The axiom does not say that Nature maximizes stability, minimizes vacuum energy, or
places every effective theory exactly at a transition.  It applies only when the
competing phases are independently admitted and no prior public fact distinguishes
them.  It then forbids an unrecorded relative bias.  Mathematically its simplest
realization is the multiple-point criticality condition studied by Froggatt and
Nielsen \cite{froggatt-nielsen-mpp}; the motivation and jurisdiction here are the
Two-Law record and provenance clauses, not their proposed microscopic mechanism.
\end{remark}

\subsection{Higgs criticality conditions}

\begin{theorem}[Two stationary degenerate Higgs phases]
\label{thm:cba-higgs-double-zero}
Suppose one unbiased phase is the electroweak vacuum and a second lies at a
nonzero high field $h_*$ for which the RG-improved effective potential has the
asymptotic form
\[
 V_{\mathrm{eff}}(h)
 =\frac14\lambda_{\mathrm{eff}}(h)h^4,
 \qquad h\gg v.
\]
After subtracting their common additive vacuum value, Critical Boundary
Actualization implies
\[
 \boxed{
 \lambda_{\mathrm{eff}}(h_*)=0,
 \qquad
 \beta_{\lambda_{\mathrm{eff}}}(h_*)=0.}
\]
At leading order in an $\overline{\mathrm{MS}}$ matching treatment these become
$\lambda(\mu_*)=\beta_\lambda(\mu_*)=0$ with $\mu_*\simeq h_*$.
\end{theorem}

\begin{proofsketch}
Degeneracy with the subtracted electroweak vacuum gives
$V_{\mathrm{eff}}(h_*)=0$.  Since $h_*\ne0$, this gives
$\lambda_{\mathrm{eff}}(h_*)=0$.  Differentiating with respect to $h$ gives
\[
 \frac{dV_{\mathrm{eff}}}{dh}
 =h^3\left(\lambda_{\mathrm{eff}}
       +\frac14\beta_{\lambda_{\mathrm{eff}}}\right).
\]
Stationarity at $h_*$ then gives the beta-function condition.  The value of the
full effective potential at an extremum is the physical statement; identifying it
with a running $\overline{\mathrm{MS}}$ quartic is a controlled matching
approximation \cite{buttazzo-criticality}.
\end{proofsketch}

\begin{theorem}[Leading Standard Model critical sum rule]
\label{thm:cba-sm-sum-rule}
Use $g_1=\sqrt{5/3}\,g_Y$ and retain the top Yukawa as the dominant Yukawa.
At one loop,
\[
\begin{aligned}
16\pi^2\beta_\lambda
={}&24\lambda^2-6y_t^4
 +\left(12y_t^2-9g_2^2-\frac95g_1^2\right)\lambda\\
&+\frac{27}{200}g_1^4
 +\frac9{20}g_1^2g_2^2
 +\frac98g_2^4.
\end{aligned}
\]
Consequently the double-zero condition gives
\[
 \boxed{
 y_t^4(\mu_*)
 =\frac16\left(
 \frac{27}{200}g_1^4
 +\frac9{20}g_1^2g_2^2
 +\frac98g_2^4
 \right)_{\mu_*}.}
\]
This fixes one dimensionless Yukawa combination from the gauge couplings rather
than inserting $y_t$ independently.
\end{theorem}

\begin{corollary}[Exact equal-trace critical ratio]
\label{cor:cba-equal-trace-ratio}
If the critical placement also obeys the already derived electroweak equal-trace
normalization $g_1(\mu_*)=g_2(\mu_*)=:g_X$, then
\[
 \boxed{
 \frac{y_t}{g_X}
 =\left(\frac{57}{200}\right)^{1/4}
 =0.7306530727\ldots .}
\]
In a vacuum in which both masses are generated by the same Higgs value,
\[
 \boxed{
 \frac{m_t}{m_W}
 =\sqrt2\frac{y_t}{g_2}
 =1.033299485\ldots }
\]
at that same running scale.  This is a scale-local running-mass relation, not the
ratio of present pole masses.
\end{corollary}

\subsection{Placement tests rather than post-hoc fitting}

\begin{proposition}[The electroweak equal-trace crossing is not the critical scale]
\label{prop:cba-crossing-killed}
The existing two-loop audit gives
\[
 \mu_X=1.08939\times10^{13}\ \mathrm{GeV},
 \qquad
 g_1(\mu_X)=g_2(\mu_X)=0.545701,
 \qquad
 g_3(\mu_X)=0.580413
\]
\cite{lande-hopf-audit}.  Imposing the critical boundary there gives
\[
 y_t(\mu_X)=0.3987181124,
 \qquad
 \lambda(\mu_X)=0.
\]
A one-loop Standard Model run to $173$ GeV gives
\[
\begin{array}{c|ccccc}
 &g_1&g_2&g_3&y_t&\lambda\\
\hline
173\ \mathrm{GeV}
&0.463774&0.650835&1.144141&0.814431&0.071325
\end{array}
\]
and the tree-level diagnostic values
\[
 m_t^{\overline{\mathrm{MS}}}\simeq141.80\ \mathrm{GeV},
 \qquad
 m_h\simeq93.00\ \mathrm{GeV}.
\]
The mismatch is much larger than a credible loop or scheme correction.  Therefore
\[
 \boxed{\mu_*\ne\mu_X}
\]
in the minimal Standard Model realization.  The attractive identification of the
structural $g_1=g_2$ crossing with Higgs phase criticality is falsified rather than
retained as numerology.
\end{proposition}

\begin{proofsketch}
Integrate
\[
16\pi^2\beta_{g_i}=b_i g_i^3,
\quad
(b_1,b_2,b_3)=\left(\frac{41}{10},-\frac{19}{6},-7\right),
\]
\[
16\pi^2\beta_{y_t}
=y_t\left(\frac92y_t^2-\frac{17}{20}g_1^2-\frac94g_2^2-8g_3^2\right)
\]
together with the beta function in Theorem~\ref{thm:cba-sm-sum-rule}.  The complete
reproduction is in \path{critical_boundary_actualization_calculations.py}.
\end{proofsketch}

\begin{theorem}[Strict Planck criticality benchmark]
\label{thm:cba-planck-benchmark}
Use the NNLO weak-scale matching and three-loop Standard Model extrapolation of
Buttazzo et al. \cite{buttazzo-criticality}.  At their central Planck-scale gauge
values
\[
 g_1(M_{\mathrm{Pl}})=0.6154,
 \qquad
 g_2(M_{\mathrm{Pl}})=0.5055,
\]
the leading critical sum rule gives
\[
 \boxed{y_t^{\mathrm{crit}}(M_{\mathrm{Pl}})=0.3882760941.}
\]
At the reference $\alpha_s(M_Z)=0.1184$, inverting their interpolation formulae
for $y_t(M_{\mathrm{Pl}})$ and $\lambda(M_{\mathrm{Pl}})$ with
$\lambda(M_{\mathrm{Pl}})=0$ gives
\[
 \boxed{
 M_t^{\mathrm{crit}}=174.47\ \mathrm{GeV},
 \qquad
 M_h^{\mathrm{crit}}=132.66\ \mathrm{GeV},}
\]
and therefore
\[
 \boxed{\frac{M_h}{M_t}=0.760341.}
\]
These numbers define the strict no-threshold, pure-Standard-Model realization of
the axiom at the Planck scale.
\end{theorem}

\begin{proofsketch}
The critical Yukawa is the fourth root in
Theorem~\ref{thm:cba-sm-sum-rule}.  The high-order interpolation used is
\[
 y_t(M_{\mathrm{Pl}})
 =0.3825+0.0051(M_t-173.34)
 -0.0021\frac{\alpha_s-0.1184}{0.0007},
\]
\[
\begin{aligned}
 \lambda(M_{\mathrm{Pl}})
={}&-0.0143-0.0066(M_t-173.34)
 +0.0018\frac{\alpha_s-0.1184}{0.0007}\\
&+0.0029(M_h-125.15),
\end{aligned}
\]
with masses in GeV.  Solve the first equation for $M_t$ and the second for $M_h$.
The calculation is algebraic once the gauge boundary is specified.
\end{proofsketch}

\begin{proposition}[Empirical verdict on the strict realization]
\label{prop:cba-strict-verdict}
The 2026 Particle Data Group summary gives
\[
 M_h=125.13\pm0.11\ \mathrm{GeV},
\]
\[
 M_t^{\mathrm{direct}}=172.60\pm0.27\ \mathrm{GeV},
 \qquad
 M_t^{\mathrm{pole\ from\ cross\ section}}=172.1\pm0.6\ \mathrm{GeV}
\]
\cite{pdg2026}.  The strict benchmark is high by $7.53$ GeV in the Higgs mass and
by $1.87$ GeV relative to the direct top average; its predicted mass ratio
$0.76034$ compares with $0.72497$ using the direct top value and $0.72708$ using
the cross-section pole extraction.  The top comparison carries the familiar
mass-definition ambiguity, but the Higgs discrepancy is decisive.  Hence:
\[
 \boxed{
 \begin{gathered}
 \text{Critical Boundary Actualization is not realized by the}\\
 \text{threshold-free minimal Standard Model at }M_{\mathrm{Pl}}.
 \end{gathered}}
\]
The axiom, if retained, predicts nonzero ultraviolet matching corrections rather
than the measured masses directly from the minimal Standard Model.
\end{proposition}

\subsection{A quantitative ultraviolet target}

\begin{theorem}[Required critical matching corrections]
\label{thm:cba-threshold-target}
Write the full ultraviolet matching at the certified gravity interface as
\[
 \lambda_{\mathrm{full}}
 =\lambda_{\mathrm{SM}}+\Delta\lambda_*,
 \qquad
 y_{t,\mathrm{full}}
 =y_{t,\mathrm{SM}}+\Delta y_{t,*},
\]
and, for this benchmark, retain the minimal-Standard-Model gauge extrapolation.
At $\alpha_s(M_Z)=0.1184$, the PDG direct top average and Higgs mass imply
\[
 \lambda_{\mathrm{SM}}(M_{\mathrm{Pl}})=-0.009474,
 \qquad
 y_{t,\mathrm{SM}}(M_{\mathrm{Pl}})=0.378726.
\]
Criticality therefore requires
\[
 \boxed{
 \Delta\lambda_*=+0.009474,
 \qquad
 \Delta y_{t,*}=+0.009550,}
\]
with the Yukawa correction equal to $2.52\%$ of the extrapolated value.
Using instead the PDG cross-section pole value gives
\[
 \boxed{
 \Delta\lambda_*=+0.006174,
 \qquad
 \Delta y_{t,*}=+0.012100,}
\]
a $3.22\%$ Yukawa correction.

For
\[
 \Delta_\alpha
 :=\frac{\alpha_s(M_Z)-0.1184}{0.0007},
\]
the direct-top targets become
\[
 \Delta\lambda_*
 =0.009474-0.0018\Delta_\alpha,
 \qquad
 \Delta y_{t,*}
 =0.009550+0.0021\Delta_\alpha.
\]
These signs and magnitudes are numerical predictions for a future explicit
quantum-gravity or ultraviolet matching calculation, conditional on the axiom and
on negligible gauge-coupling threshold shifts.
\end{theorem}

\begin{corollary}[Observed location relative to the stability boundary]
At the reference strong coupling, the high-order Standard Model stability line is
\[
 M_h^{\mathrm{stab}}
 \simeq129.6\ \mathrm{GeV}+2(M_t-173.34\ \mathrm{GeV})
 \quad(\pm0.3\ \mathrm{GeV\ perturbative})
\]
\cite{buttazzo-criticality}.  It gives
\[
 M_h^{\mathrm{stab}}=128.12\ \mathrm{GeV}
\]
for the direct top average and $127.12$ GeV for the cross-section pole extraction.
The measured $125.13$ GeV lies respectively $2.99$ GeV and $1.99$ GeV below these
central boundaries.  Thus the observed minimal-Standard-Model trajectory lies on
the metastable side but close enough that an order-$10^{-2}$ ultraviolet quartic
matching correction would move it to exact criticality.
\end{corollary}

\begin{remark}[What has and has not been predicted]
The new axiom produces genuine numerical content:
\[
 \lambda_*=0,
 \qquad
 y_t^4=\frac16\left(
 \frac{27}{200}g_1^4+\frac9{20}g_1^2g_2^2+\frac98g_2^4\right),
\]
the exact equal-trace ratio $(57/200)^{1/4}$, a strict Planck benchmark, a killed
$10^{13}$ GeV placement, and required ultraviolet threshold corrections.
It does not yet determine the lighter Yukawa hierarchy, CKM or PMNS angles, the
Koide high root, the absolute electroweak scale, Newton's constant, or the common
additive cosmological constant.  Vacuum degeneracy fixes differences of vacuum
functionals; it does not fix their common additive value.  Those remain separate
provenance problems.
\end{remark}

\begin{remark}[Closure ledger after Critical Boundary Actualization]
\begin{center}
\small
\begin{tabularx}{\textwidth}{l Y}
\toprule
Target & Status \\
\midrule
Higgs quartic at the critical interface
& $\lambda_{\mathrm{eff}}=0$ and $\beta_{\lambda_{\mathrm{eff}}}=0$ derived. \\
Top Yukawa combination
& Fixed by the gauge-quartic critical sum rule. \\
Equal-trace dimensionless ratio
& $y_t/g_X=(57/200)^{1/4}=0.730653$ and $m_t/m_W=1.033299$ at the same critical scale. \\
Critical-scale placement at $g_1=g_2$
& Refuted by the one-loop low-energy diagnostic; $\mu_*\ne1.09\times10^{13}$ GeV in the minimal SM. \\
Threshold-free Planck realization
& Predicts $(M_t,M_h)=(174.47,132.66)$ GeV and fails the measured Higgs mass. \\
Threshold-aware Planck realization
& Requires positive matching corrections $\Delta\lambda_*\simeq0.0062$--$0.0095$ and $\Delta y_{t,*}\simeq0.0096$--$0.0121$ at the reference $\alpha_s$. \\
Mass ratio
& Strict benchmark $M_h/M_t=0.76034$; measured ratio is about $0.725$--$0.727$. \\
Vacuum status
& Minimal-SM trajectory lies on the metastable side of the stability boundary. \\
Lighter masses and mixing
& Not fixed; criticality gives a top-dominated trace relation rather than the flavor spectrum. \\
Cosmological constant
& Common additive vacuum value remains undetermined. \\
\bottomrule
\end{tabularx}
\end{center}
\end{remark}
